% !TeX root = ../sections/02_exercises.tex

\begin{exercise}
	R-5. $A$ を有限生成自由アーベル群とし，$B$ をその部分群とする。
	次を示せ。
	\begin{enumerate}
		\item $B$ の階数は $A$ の階数以下である。
		\item $A/B$ が有限群であることと，
		\[
		\operatorname{rank} A=\operatorname{rank} B
		\]
		であることは同値である。
	\end{enumerate}
\end{exercise}

\begin{background}
	有限生成自由アーベル群とは，
	ある非負整数 $n$ に対して
	\[
	A\cong \mathbb{Z}^n
	\]
	と表されるアーベル群のことである。
	この $n$ を $A$ の階数といい，
	\[
	\operatorname{rank}A=n
	\]
	と書く。
	
	自由アーベル群の部分群は再び自由アーベル群である。
	したがって，$B\subset A$ も有限生成自由アーベル群であり，
	$\operatorname{rank}B$ が定義できる。
	
	階数は，有理数体 $\mathbb{Q}$ でテンソルすることで
	ベクトル空間の次元として捉えることができる。
	すなわち，
	\[
	\operatorname{rank}A
	=
	\dim_{\mathbb{Q}}(A\otimes_{\mathbb{Z}}\mathbb{Q})
	\]
	である。
	
	また，アーベル群 $G$ が有限群であるならば，
	\[
	G\otimes_{\mathbb{Z}}\mathbb{Q}=0
	\]
	である。
	逆に，有限生成アーベル群 $G$ について
	\[
	G\otimes_{\mathbb{Z}}\mathbb{Q}=0
	\]
	であれば，$G$ は階数 $0$ であり，したがって有限群である。
\end{background}

\begin{strategy}
	包含写像
	\[
	B\hookrightarrow A
	\]
	を考える。
	これに $\mathbb{Q}$ とのテンソル積を取ると，
	\[
	B\otimes_{\mathbb{Z}}\mathbb{Q}
	\longrightarrow
	A\otimes_{\mathbb{Z}}\mathbb{Q}
	\]
	という $\mathbb{Q}$-ベクトル空間の単射が得られる。
	
	したがって，
	\[
	\dim_{\mathbb{Q}}(B\otimes_{\mathbb{Z}}\mathbb{Q})
	\leq
	\dim_{\mathbb{Q}}(A\otimes_{\mathbb{Z}}\mathbb{Q})
	\]
	である。
	これより
	\[
	\operatorname{rank}B\leq \operatorname{rank}A
	\]
	が従う。
	
	次に，完全列
	\[
	0\longrightarrow B\longrightarrow A\longrightarrow A/B\longrightarrow 0
	\]
	を考える。
	これに $\mathbb{Q}$ とのテンソル積を取ることで，
	階数について
	\[
	\operatorname{rank}(A/B)
	=
	\operatorname{rank}A-\operatorname{rank}B
	\]
	が得られる。
	
	したがって，
	\[
	\operatorname{rank}A=\operatorname{rank}B
	\]
	であることは，
	\[
	\operatorname{rank}(A/B)=0
	\]
	であることと同値である。
	
	$A/B$ は有限生成アーベル群であるから，階数が $0$ であることは
	$A/B$ が有限群であることと同値である。
	
	よって，
	\[
	A/B\text{ が有限群}
	\quad\Longleftrightarrow\quad
	\operatorname{rank}A=\operatorname{rank}B
	\]
	が成り立つ。
\end{strategy}